% Compile from the repository root:
% lualatex --interaction=nonstopmode --halt-on-error --output-directory=data test-results.tex
\documentclass[fontsize=10pt,paper=b5,twoside,DIV=11,BCOR=9mm,
    headinclude=true,footinclude=false]{scrartcl}

% Font setup from dissertation/src/packages.tex and style/presets.tex.
\usepackage[libertine]{newtxmath}
\usepackage[no-math]{fontspec}
\setmainfont{Linux Libertine O}
\setsansfont{Linux Biolinum O}
\setmonofont[Scale=0.875]{Droid Sans Mono}
\usepackage{amsmath}
\usepackage[table]{xcolor}
\usepackage{array,booktabs,multirow,ifthen,subcaption}
\usepackage{pgfplots}
\pgfplotsset{compat=1.18}
\usepackage{currfile,xurl}

% Label each plot with its source filename, keeping the label and plot together.
% These hooks apply only to this preview; individual plot files stay unchanged.
\AddToHook{env/tikzpicture/before}{%
    \par\noindent\begin{minipage}{\linewidth}\raggedright
    {\scriptsize\expandafter\url\expandafter{\currfilename}\par}\smallskip
}
\AddToHook{env/tikzpicture/after}{\end{minipage}\par}

% Only the symbols and macros needed by the generated tables, copied from
% dissertation/src/style/customsymbols.tex and dissertation/src/commands.tex.
\makeatletter
\DeclareFontEncoding{LS2}{}{\@noaccents}
\makeatother
\DeclareFontSubstitution{LS2}{stix}{m}{n}
\DeclareSymbolFont{largesymbolsstix}{LS2}{stixex}{m}{n}
\DeclareMathDelimiter{\lbrbrak}{\mathopen}{largesymbolsstix}{"EE}{largesymbolsstix}{"14}
\DeclareMathDelimiter{\rbrbrak}{\mathclose}{largesymbolsstix}{"EF}{largesymbolsstix}{"15}
\newcommand*{\tuplestyle}[1]{\boldsymbol{#1}}
\newcommand*{\tupleaccess}[2]{{#1\!\lbrbrak#2\rbrbrak}}
\newcommand*{\vecrew}{\tuplestyle{R}}
\newcommand*{\boundtuple}{\tuplestyle{t}}
\newcommand*{\costthresh}{\boundtuple} % Name used in the original plot formulation.
\newcommand{\leveldegree}{\tuplestyle{d}}
\newcommand*{\model}[1]{\textsf{#1}}
\newcommand*{\sizethreshold}{\tau}
\newcommand*{\resolution}{\eta}
\newcommand*{\config}[2][]{\ifthenelse{\equal{#1}{}}{\textsc{#2}}{\textsc{#2}$_{#1}$}}

\setlength{\parindent}{0pt}
\raggedbottom

\begin{document}

\textbf{Benchmark instances}\par\medskip
{\small\input{data/rewbnd-latex/tablebenchmarks.tex}\par}

\clearpage
\textbf{Unfold results}\par\medskip
\input{data/rewbnd-latex/tableunf.tex}

\clearpage
\textbf{Reward-aware results (1800 seconds)}\par\medskip
\input{data/rewbnd-latex/tablet1800.tex}

% The collection indices below input one separate file per plot.
% Each collection starts on a fresh page. To test an individual plot, use
% its filename instead, e.g. data/rew-bounded-plots/bnds_result-clean.tex.
\clearpage
\setlength{\parskip}{1em}
\section*{Time versus value bounds}
\input{data/rew-bounded-plots/time_result.tex}

\clearpage
\section*{Shared legend without externalization}
\begin{figure}[!ht]
  \centering
  \begin{subfigure}[t]{0.49\linewidth}
    \centering
    \input{data/rew-bounded-plots/nolegend/lvls_result-clean-nolegend.tex}
    \caption{Levels versus value bounds.}
  \end{subfigure}
  \hfill
  \begin{subfigure}[t]{0.49\linewidth}
    \centering
    \input{data/rew-bounded-plots/nolegend/bnds_result-clean-nolegend.tex}
    \caption{Reward bounds versus value bounds.}
  \end{subfigure}

  \medskip
  \makebox[\linewidth][c]{\input{data/rew-bounded-plots/lvls_result-legend.tex}}
  \caption{Two generated plots with one shared method legend.}
\end{figure}

\end{document}
